\chapter{Blood Culture}

\section*{What Is This Test?}
A blood culture detects bacteria or fungi circulating in the bloodstream. Blood is placed into specialized aerobic and anaerobic culture bottles and monitored for microbial growth, after which the organism is identified and antimicrobial susceptibility testing is performed.

\section*{Alternative Names}
Blood Culture and Sensitivity; Blood C\&S; Bacteremia Culture; Sepsis Culture.

\section*{Why This Test Is Ordered}
Blood cultures are essential when bloodstream infection, sepsis, infective endocarditis, meningitis with systemic illness, or an unexplained fever with high-risk features is suspected. They guide targeted antimicrobial treatment and can identify resistant organisms.

\section*{How This Test Is Performed}
Usually two or more culture sets are collected from separate venipuncture sites, ideally before antibiotics when this is clinically safe. Each set commonly includes one aerobic and one anaerobic bottle. The laboratory monitors bottles for growth, performs a Gram stain when positive, and identifies the organism with susceptibility testing.

\section*{Preparation, Risks, and Limitations}
Fasting is not required. Careful skin antisepsis and adequate blood volume are critical. Risks are those of venipuncture. Antibiotics given before collection can reduce recovery, while skin contamination can produce a misleading positive result. A single bottle growing common skin flora may represent contamination, but interpretation depends on the organism, number of positive bottles, timing, and clinical context.

\section*{Understanding Results}
No growth at the final reporting time is reassuring but does not exclude infection, especially after antibiotics or with fastidious organisms. A positive culture is reported urgently; the organism and susceptibility profile determine whether the result represents true bacteremia, contamination, or a polymicrobial infection.

\section*{References}
Clinical and Laboratory Standards Institute guidance for blood culture collection and microbiology reporting.

\clearpage
\section*{Understanding Results and Follow-up}
The laboratory report should identify the specimen, method, units, reference interval or decision limit, and whether the result is positive, negative, abnormal, or inconclusive. Reference intervals vary with age, sex, pregnancy, specimen type, assay, and laboratory; a result outside a range is not automatically a diagnosis, and a result within a range does not always exclude disease. Discuss unexpected results with the ordering clinician, who can decide whether repeat or confirmatory testing, treatment, or referral is appropriate. Do not change medicines or delay urgent care based on an isolated result.


\section*{Regulatory Status and Authorization Date}
This chapter describes a test category, not a specific commercial product. FDA, CDSCO, EU IVDR, and MHRA approval or registration dates therefore cannot be assigned without the assay manufacturer, product name, instrument, and intended use. For laboratory-developed tests, an individual agency approval date may not exist. See the Preface for the interpretation of regulatory status.

\clearpage
